% ================================================================
% SECTION 9: STATIC CONDENSATION
% ================================================================
\section{Static Condensation}
\label{sec:static_condensation}

\begin{center}
\code{bionetflux.core.static\_condensation\_base} ~~\textbar~~
\code{bionetflux.core.static\_condensation\_factory} \\
\code{bionetflux.core.static\_condensation\_keller\_segel} ~~\textbar~~
\code{bionetflux.core.static\_condensation\_ooc}
\end{center}

Static condensation eliminates the interior (bulk and flux) unknowns
element-by-element, reducing the global system to trace unknowns only.
The module uses an ABC $\to$ Factory $\to$ Concrete-class hierarchy.

% ------------------------------------------------------------------
\subsection{StaticCondensationBase (ABC)}
\label{sec:sc_base}

\begin{lstlisting}[language=Python, caption={Base class constructor}]
class StaticCondensationBase(ABC):
    def __init__(self, problem: Problem,
                 global_disc: GlobalDiscretization,
                 elementary_matrices: Any, ipb: int = 0):
\end{lstlisting}

\noindent Stored attributes:

\begin{longtable}{p{5cm}p{9.5cm}}
\toprule
\textbf{Attribute} & \textbf{Description} \\
\midrule
\code{problem} & Reference to the \code{Problem} object \\
\code{discretization} & \code{global\_disc.spatial\_discretizations[ipb]} \\
\code{elementary\_matrices} & Shared \code{ElementaryMatrices} instance \\
\code{sc\_matrices} & Dictionary filled by \code{build\_matrices()} \\
\code{dt} & Time step, sourced from \code{global\_disc.dt} \\
\code{tau} & Stabilization vector from \code{discretization.tau} \\
\code{flux\_orders} & Empty list---must be set by the subclass \\
\bottomrule
\end{longtable}

\subsubsection{Abstract Methods}

Every concrete subclass must implement:

\begin{lstlisting}[language=Python, caption={Abstract interface}]
@abstractmethod
def build_matrices(self) -> Dict[str, np.ndarray]: ...

@abstractmethod
def static_condensation(self, local_trace: np.ndarray,
                        local_source: np.ndarray,
                        **kwargs) -> Tuple[np.ndarray, np.ndarray,
                                           np.ndarray, np.ndarray]:
    """Returns (local_solution, flux, flux_trace, jacobian)."""

@abstractmethod
def assemble_forcing_term(self, *args, **kwargs) -> np.ndarray: ...
\end{lstlisting}

\subsubsection{Flux-Order Machinery}

\begin{structurebox}{Flux orders and DOF counts}
\code{flux\_orders} is a list of length $n_\mathrm{eq}$.  Entry $k$:
\begin{itemize}
  \item \textbf{0} $\Rightarrow$ P0 flux $\Rightarrow$ 1 DOF per element.
  \item \textbf{1} $\Rightarrow$ P1 flux $\Rightarrow$ 2 DOFs per element.
\end{itemize}
The base class provides:
\end{structurebox}

\begin{longtable}{p{7cm}p{7.5cm}}
\toprule
\textbf{Property / Method} & \textbf{Description} \\
\midrule
\code{flux\_dofs\_per\_element -> List[int]} & \code{[order + 1 for order in flux\_orders]} \\
\code{total\_flux\_dofs\_per\_element -> int} & Sum of the above \\
\code{\_validate\_flux\_orders()} & Checks length = \code{neq} and values $\in \{0,1\}$ \\
\bottomrule
\end{longtable}

\subsubsection{update\_dt}

\begin{lstlisting}[language=Python, caption={Adaptive time-step update}]
def update_dt(self) -> None:
    """Re-read dt from GlobalDiscretization and rebuild matrices."""
    new_dt = self._global_disc.dt
    self.dt = new_dt
    self.build_matrices()
\end{lstlisting}

Called by \code{AdaptiveTimeStepper} after adjusting the global $\Delta t$.

% ------------------------------------------------------------------
\subsection{StaticCondensationFactory}

\begin{lstlisting}[language=Python, caption={Factory class}]
class StaticCondensationFactory:
    _implementations = {
        "keller_segel": KellerSegelStaticCondensation,
        "organ_on_chip": StaticCondensationOOC,
    }

    @classmethod
    def create(cls, problem, global_disc, elementary_matrices,
               i: int = 0) -> StaticCondensationBase: ...

    @classmethod
    def register_implementation(cls, problem_type: str,
                                implementation_class): ...

    @classmethod
    def get_available_types(cls) -> list: ...
\end{lstlisting}

The factory dispatches on \code{problem.type} (a string).  New problem types
can be registered at runtime via \code{register\_implementation}.

% ------------------------------------------------------------------
\subsection{KellerSegelStaticCondensation}
\label{sec:sc_ks}

\begin{lstlisting}[language=Python, caption={Keller-Segel condensation}]
class KellerSegelStaticCondensation(StaticCondensationBase):
    def __init__(self, problem, global_disc, elementary_matrices, ipb=0):
        super().__init__(problem, global_disc, elementary_matrices, ipb)
        self.flux_orders = [0, 1]   # P0 for u, P1 for phi
\end{lstlisting}

\noindent The two-equation KS system:
\begin{itemize}
  \item Equation 0 ($u$): P0 flux (1 DOF per element).
  \item Equation 1 ($\varphi$): P1 flux (2 DOFs per element).
\end{itemize}

\code{build\_matrices()} extracts parameters $\mu$, $\nu$, $a$, $b$ from
\code{problem.parameters[0:4]} and computes HDG element matrices using
the elementary matrices \code{M}, \code{D}, \code{T}, \code{IM}, \code{Mb}, \code{Gb}, \code{Ntil}, \code{Nhat}, \code{QUAD}.
\code{static\_condensation(local\_trace, local\_source)} returns
\code{(local\_solution, flux, flux\_trace, jacobian)} for one element.

% ------------------------------------------------------------------
\subsection{StaticCondensationOOC}
\label{sec:sc_ooc}

\begin{lstlisting}[language=Python, caption={Organ-on-Chip condensation}]
class StaticCondensationOOC(StaticCondensationBase):
    def __init__(self, problem, global_disc, elementary_matrices, ipb=0):
        super().__init__(problem, global_disc, elementary_matrices, ipb)
        self.flux_orders = [0, 1, 1, 1]  # P0 for u, P1 for w, v, phi
\end{lstlisting}

\noindent The four-equation OoC system:
\begin{itemize}
  \item Equation 0 ($u$):  P0 flux (1 DOF per element).
  \item Equation 1 ($\omega$): P1 flux (2 DOFs per element).
  \item Equation 2 ($v$):  P1 flux (2 DOFs per element).
  \item Equation 3 ($\varphi$): P1 flux (2 DOFs per element).
\end{itemize}

\code{build\_matrices()} extracts nine parameters $\nu, \mu, \epsilon, \sigma, a, b, c, d, \chi$ from
\code{problem.parameters[0:9]}.  The internal structure mirrors the Keller-Segel implementation scaled to four equations.
